\subsection{Isometries and Unitary Matrices}

Throughout this subsection, $\mathbb{F}\in\{\mathbb{R},\mathbb{C}\}$.

\begin{definition} \label{def:isometry-linear-map} \hlt{Isometry and Isometric Isomorphism}\\
Let $\mathcal{V}$ and $\mathcal{W}$ be inner-product spaces and let $T\in\LL(\mathcal{V},\mathcal{W})$.
\begin{enumerate}[label=(\roman*)]
\setlength{\itemsep}{0pt}
\item $T$ is an \hlt{isometry} if $(\forall\,\mathbf{v}\in\mathcal{V})\,\norm{T\mathbf{v}}=\norm{\mathbf{v}}$.
\item $T$ is an \hlt{isometric isomorphism} if it is a bijective isometry. Spaces $\mathcal{V}$ and $\mathcal{W}$ are \hlt{isometrically isomorphic} if such a map $T$ exists.
\end{enumerate}
\end{definition}

\begin{proposition} \label{prop:basic-isometry-properties} \hlt{Basic Properties of Isometries}\\
Let $\mathcal{V}$ be an inner-product space and let $S,T\in\LL(\mathcal{V})$ be isometries.
\begin{enumerate}[label=(\roman*)]
\setlength{\itemsep}{0pt}
\item $ST$ is an isometry;
\item $T$ is injective;
\item if $\mathcal{V}$ is finite-dimensional, then $T$ is invertible;
\item if $\mathcal{V}$ is finite-dimensional, then $T^{-1}$ is an isometry.
\end{enumerate}
\begin{proof}
\begin{enumerate}[label=(\roman*)]
\setlength{\itemsep}{0pt}
\item $(\forall\,\mathbf{v}\in\mathcal{V})\,\norm{ST\mathbf{v}}=\norm{T\mathbf{v}}=\norm{\mathbf{v}}$.
\item If $\mathbf{v}\in\ker(T)$, then $\norm{\mathbf{v}}=\norm{T\mathbf{v}}=0$, so $\mathbf{v}=\mathbf{0}$ by positive definiteness. Hence $\ker(T)=\{\mathbf{0}\}$, and $T$ is injective.
\item $T$ is invertible by (ii) and Theorem \ref{thm:finite-dimensional-operator-invertible-equivalences}.
\item Suppose $\dim(\mathcal{V})<\infty$. Then $T^{-1}$ exists by (iii). $(\forall\,\mathbf{u}\in\mathcal{V})\,\norm{T^{-1}\mathbf{u}}=\norm{TT^{-1}\mathbf{u}}=\norm{\mathbf{u}}$.
\end{enumerate}
\hlt{Structural Commentary}: Norm identities establish closure and injectivity before finite-dimensional invertibility controls the inverse.\\
\hlt{Key Insight}: Length preservation forces the kernel to vanish.
\end{proof}
\end{proposition}

\begin{counterexample} \label{cex:l2-shift-isometry-not-surjective} \hlt{An Isometry Need Not Be Onto}\\
Define $T:\ell^2(\mathbb{F})\to\ell^2(\mathbb{F})$ by $T(x_1,x_2,\ldots)=(0,x_1,x_2,\ldots)$. The map $T$ is linear by coordinatewise operations.\\
For $\mathbf{x}=(x_j)_{j\geq1}$, $\norm{T\mathbf{x}}_2^2=\sum\limits_{j=1}^{\infty}\abs{x_j}^2=\norm{\mathbf{x}}_2^2$, so $T$ is an isometry.\\
Every vector in $\operatorname{range}(T)$ has first coordinate $0$, whereas $(1,0,0,\ldots)\notin\operatorname{range}(T)$.\\
Thus $T$ is not surjective, so an infinite-dimensional isometry need not be an isomorphism.
\end{counterexample}

\begin{theorem} \label{thm:isometry-iff-inner-product-preserving} \hlt{Isometries Preserve Inner Products}\\
Let $\mathcal{V}$ and $\mathcal{W}$ be inner-product spaces and let $T\in\LL(\mathcal{V},\mathcal{W})$. Then $(\forall\,\mathbf{x}\in\mathcal{V})\,\norm{T\mathbf{x}}=\norm{\mathbf{x}}\Longleftrightarrow(\forall\,\mathbf{u},\mathbf{v}\in\mathcal{V})\,\innerprod[\mathcal{W}]{T\mathbf{u},T\mathbf{v}}=\innerprod[\mathcal{V}]{\mathbf{u},\mathbf{v}}$.
\begin{proof}
\begin{enumerate}[label=\roman*.]
\setlength{\itemsep}{0pt}
\item[$(\Rightarrow)$] Suppose $T$ preserves norms and fix $\mathbf{u},\mathbf{v}\in\mathcal{V}$.\\
By linearity and norm preservation, the polarization formulas for $(T\mathbf{u},T\mathbf{v})$ and $(\mathbf{u},\mathbf{v})$ have equal norm terms.\\
Thus $\innerprod[\mathcal{W}]{T\mathbf{u},T\mathbf{v}}=\innerprod[\mathcal{V}]{\mathbf{u},\mathbf{v}}$ by Theorem \ref{thm:polarization-identities}.
\item[$(\Leftarrow)$] If $T$ preserves inner products, then $(\forall\,\mathbf{u}\in\mathcal{V})\,\norm{T\mathbf{u}}^2=\innerprod[\mathcal{W}]{T\mathbf{u},T\mathbf{u}}=\innerprod[\mathcal{V}]{\mathbf{u},\mathbf{u}}=\norm{\mathbf{u}}^2$.\\
Thus $T$ preserves norms by nonnegativity.
\end{enumerate}
\hlt{Structural Commentary}: Polarization transfers the norm identity to mixed inner products; diagonal evaluation reverses the transfer.\\
\hlt{Key Insight}: A norm determines its inducing inner product.
\end{proof}
\end{theorem}

\begin{remark} \label{rem:isometry-conventions} \hlt{Isometry Conventions}\\
For linear maps, norm and inner-product preservation are equivalent by Theorem \ref{thm:isometry-iff-inner-product-preserving}.\\
Thus either convention defines the same isometries.\\
For possibly degenerate bilinear forms, form preservation need not imply injectivity.\\
An isometry is therefore often required to be bijective and form-preserving.
\end{remark}

\begin{theorem} \label{thm:finite-dimensional-isometry-iff-adjoint-product-identity} \hlt{Isometries and the Adjoint Product}\\
Let $\mathcal{V}$ be a finite-dimensional inner-product space and let $T\in\LL(\mathcal{V})$. Then $(\forall\,\mathbf{v}\in\mathcal{V})\,\norm{T\mathbf{v}}=\norm{\mathbf{v}}\Longleftrightarrow T^*T=I_{\mathcal{V}}$.
\begin{proof}
\begin{enumerate}[label=\roman*.]
\setlength{\itemsep}{0pt}
\item[$(\Rightarrow)$] Suppose $T$ preserves norms and set $A\defeq T^*T-I_{\mathcal{V}}$. Then $(\forall\,\mathbf{u},\mathbf{v}\in\mathcal{V})\,\innerprod{\mathbf{u},A\mathbf{v}}=0$ by adjointness and Theorem \ref{thm:isometry-iff-inner-product-preserving}.\\
Taking $\mathbf{u}=A\mathbf{v}$, $(\forall\,\mathbf{v}\in\mathcal{V})\,\norm{A\mathbf{v}}^2=0$. Hence $A=0$ by positive definiteness, so $T^*T=I_{\mathcal{V}}$.
\item[$(\Leftarrow)$] If $T^*T=I_{\mathcal{V}}$, then $(\forall\,\mathbf{u}\in\mathcal{V})\,\norm{T\mathbf{u}}^2=\innerprod{T\mathbf{u},T\mathbf{u}}=\innerprod{\mathbf{u},T^*T\mathbf{u}}=\norm{\mathbf{u}}^2$ by the adjoint relation.\\
Thus $T$ preserves norms by nonnegativity.
\end{enumerate}
\hlt{Structural Commentary}: Adjointness converts preserved inner products into the vanishing of $T^*T-I_{\mathcal{V}}$.\\
\hlt{Key Insight}: The Gram operator of an isometry is the identity.
\end{proof}
\end{theorem}

\begin{definition} \label{def:unitary-operator-and-matrix} \hlt{Unitary and Orthogonal}
\begin{enumerate}[label=(\roman*)]
\setlength{\itemsep}{0pt}
\item Let $\mathcal{V}$ be a finite-dimensional inner-product space. An operator $U\in\LL(\mathcal{V})$ is \hlt{unitary} if $U^*U=UU^*=I_{\mathcal{V}}$.
\item Let $n\geq1$. A matrix $U\in\operatorname{Mat}(n,n,\mathbb{F})$ is \hlt{unitary} if $U^*U=I_n$. If $\mathbb{F}=\mathbb{R}$, a unitary matrix is \hlt{orthogonal}.
\end{enumerate}
\end{definition}

\begin{theorem} \label{thm:unitary-matrix-equivalent-conditions} \hlt{Equivalent Conditions for a Unitary Matrix}\\
Let $n\geq1$ and let $U=[\mathbf{u}_1\ \cdots\ \mathbf{u}_n]\in\operatorname{Mat}(n,n,\mathbb{F})$. The following conditions are equivalent:
\begin{enumerate}[label=(\roman*)]
\setlength{\itemsep}{0pt}
\item $U^*U=I_n$;
\item the columns $\mathbf{u}_1,\ldots,\mathbf{u}_n$ are orthonormal in $\mathbb{F}^n$;
\item $UU^*=I_n$;
\item $U^*$ is unitary;
\item the rows of $U$ are orthonormal;
\item $U$ is invertible and $U^{-1}=U^*$;
\item $(\forall\,\mathbf{x}\in\mathbb{F}^n)\,\norm{U\mathbf{x}}_2=\norm{\mathbf{x}}_2$;
\item $(\forall\,\mathbf{x},\mathbf{y}\in\mathbb{F}^n)\,\innerprod[\mathbb{F}^n]{U\mathbf{x},U\mathbf{y}}=\innerprod[\mathbb{F}^n]{\mathbf{x},\mathbf{y}}$.
\end{enumerate}
\begin{proof}
\begin{enumerate}[label=(\roman*)]
\setlength{\itemsep}{0pt}
\item[$(i)$] $\Rightarrow (ii)$ For $i,j\in J_n$, $(U^*U)_{ij}=\mathbf{u}_i^*\mathbf{u}_j=\innerprod[\mathbb{F}^n]{\mathbf{u}_j,\mathbf{u}_i}$. Hence $U^*U=I_n$ if and only if the columns are orthonormal.
\item[$(ii)$] $\Rightarrow (iii)$ Column orthonormality implies $U^*U=I_n$, so $UU^*=I_n$ by Theorem \ref{thm:square-one-sided-inverse-is-two-sided}.
\item[$(iii)$] $\Rightarrow (iv)$ The identity $(U^*)^*U^*=UU^*=I_n$ holds by conjugate-transpose involution, so $U^*$ is unitary.
\item[$(iv)$] $\Rightarrow (v)$ Unitarity of $U^*$ implies $UU^*=I_n$.\\
Let $\mathbf{r}_i$ denote row $i$ of $U$. Then $(\forall\,i,j\in J_n)\,(UU^*)_{ij}=\innerprod[\mathbb{F}^n]{\mathbf{r}_i,\mathbf{r}_j}$, so the rows are orthonormal.
\item[$(v)$] $\Rightarrow (vi)$ Row orthonormality implies $UU^*=I_n$. Thus $U$ is invertible and $U^{-1}=U^*$ by Theorem \ref{thm:square-one-sided-inverse-is-two-sided}.
\item[$(vi)$] $\Rightarrow (vii)$ The identity $U^*U=I_n$ holds by $U^{-1}=U^*$.\\
Hence $(\forall\,\mathbf{x}\in\mathbb{F}^n)\,\norm{U\mathbf{x}}_2^2=\innerprod[\mathbb{F}^n]{\mathbf{x},U^*U\mathbf{x}}=\norm{\mathbf{x}}_2^2$ by the conjugate-transpose adjoint relation.\\
Thus $U$ preserves norms by nonnegativity.
\item[$(vii)$] $\Rightarrow (viii)$ The map $\mathbf{x}\mapsto U\mathbf{x}$ preserves inner products by Theorem \ref{thm:isometry-iff-inner-product-preserving}.
\item[$(viii)$] $\Rightarrow (i)$ Set $A\defeq U^*U-I_n$. Then $(\forall\,\mathbf{x},\mathbf{y}\in\mathbb{F}^n)\,\innerprod[\mathbb{F}^n]{\mathbf{x},A\mathbf{y}}=0$ by the conjugate-transpose adjoint relation and (viii).\\
Taking $\mathbf{x}=A\mathbf{y}$, $(\forall\,\mathbf{y}\in\mathbb{F}^n)\,\norm{A\mathbf{y}}_2^2=0$. Hence $A=0$ by positive definiteness, so $U^*U=I_n$.
\end{enumerate}
\hlt{Structural Commentary}: The implication cycle connects column and row Gram identities through inverse and invariance relations.\\
\hlt{Key Insight}: Both Gram products encode orthonormal coordinate frames.
\end{proof}
\end{theorem}

Figure \ref{fig:unitary-isometry-geometry} depicts the geometric invariance in Theorem \ref{thm:unitary-matrix-equivalent-conditions}.
\begin{illustrationfigure}[H]
  \centering
  \begin{tikzpicture}[
  >=Stealth,
  every node/.style={font=\small},
  circleline/.style={gray!50, line width=0.65pt},
  veca/.style={->, blue!75!black, line width=0.95pt},
  vecb/.style={->, teal!70!black, line width=0.95pt},
  map/.style={->, gray!70, line width=0.8pt}
]
  \coordinate (O1) at (0,0);
  \coordinate (O2) at (5.0,0);
  \coordinate (Ua) at ($(O2)+(0.72,0.72)$);
  \coordinate (Ub) at ($(O2)+(-0.72,0.72)$);
  \draw[circleline] (O1) circle (1.15);
  \draw[circleline] (O2) circle (1.15);

  \draw[veca] (O1) -- (1.0,0) node[below] {$\mathbf{e}_1$};
  \draw[vecb] (O1) -- (0,1.0) node[left] {$\mathbf{e}_2$};
  \draw[veca] (O2) -- (Ua) node[above right] {$U\mathbf{e}_1$};
  \draw[vecb] (O2) -- (Ub) node[above left] {$U\mathbf{e}_2$};

  \draw[gray!55] (0.82,0) arc[start angle=0,end angle=90,radius=0.82];
  \draw[gray!55] ($(O2)+(0.58,0.58)$) arc[start angle=45,end angle=135,radius=0.82];
  \node at (0.55,0.55) {$\perp$};
  \node at ($(O2)+(0,0.78)$) {$\perp$};

  \draw[map] (1.65,0.15) -- (3.35,0.15) node[midway, above] {$U$};
\end{tikzpicture}

  \caption{A unitary map preserves the lengths and right angle of an orthonormal pair.}
  \label{fig:unitary-isometry-geometry}
\end{illustrationfigure}

\begin{definition} \label{def:unitary-and-permutation-similarity} \hlt{Unitary and Permutation Similarity}\\
Let $n\geq1$ and let $A,B\in\operatorname{Mat}(n,n,\mathbb{F})$.
\begin{enumerate}[label=(\roman*)]
\setlength{\itemsep}{0pt}
\item The matrices $A$ and $B$ are \hlt{unitarily similar} if $A=UBU^*$ for a unitary $U\in\operatorname{Mat}(n,n,\mathbb{F})$.
\item They are \hlt{orthogonally similar} if $\mathbb{F}=\mathbb{R}$ and $A=UBU^*$ for an orthogonal $U\in\operatorname{Mat}(n,n,\mathbb{R})$.
\item A matrix $P=[p_{ij}]\in\operatorname{Mat}(n,n,\mathbb{R})$ is a \hlt{permutation matrix} if $(\forall\,i,j\in J_n)\,p_{ij}\in\{0,1\}$ and $(\forall\,i\in J_n)\,\sum\limits_{j=1}^{n}p_{ij}=\sum\limits_{j=1}^{n}p_{ji}=1$.
\item The matrices $A$ and $B$ are \hlt{permutation similar} if $A=PBP^*$ for a permutation matrix $P$.
\end{enumerate}
\end{definition}

\begin{proposition} \label{prop:permutation-matrices-are-orthogonal} \hlt{Permutation Matrices Are Orthogonal}\\
Let $n\geq1$, let $\sigma\in\mathfrak{S}_n$, and set $P\defeq[\mathbf{e}_{\sigma(1)}\ \cdots\ \mathbf{e}_{\sigma(n)}]\in\operatorname{Mat}(n,n,\mathbb{R})$. Then $P^*P=PP^*=I_n$, so $P$ is orthogonal.
\begin{proof}
The columns of $P$ are a permutation of the standard orthonormal basis of $\mathbb{R}^n$.\\
Hence $P^*P=PP^*=I_n$ by Theorem \ref{thm:unitary-matrix-equivalent-conditions}(ii)--(iii). Therefore $P$ is orthogonal by the real-field definition.\\
\hlt{Structural Commentary}: Standard-basis permutation converts column orthonormality directly into both Gram identities.\\
\hlt{Key Insight}: Reordering an orthonormal basis preserves its Gram matrix.
\end{proof}
\end{proposition}

\begin{theorem} \label{thm:orthogonal-projection-matrix-unitary-coordinate-form} \hlt{Unitary Coordinate Form of an Orthogonal Projection}\\
Let $n\geq1$, let $P\in\operatorname{Mat}(n,n,\mathbb{F})$ be an orthogonal projection matrix, and set $r\defeq\operatorname{rank}(P)$.\\
Then $(\exists\,U\in\operatorname{Mat}(n,n,\mathbb{F}))\,[U^*U=I_n\land U^*PU=I_r\oplus0_{n-r}]$. Equivalently, $P$ is unitarily similar to $I_r\oplus0_{n-r}$.
\begin{proof}
\begin{enumerate}[label=(\roman*)]
\setlength{\itemsep}{0pt}
\item \hlt{Projection Operator}: Define $T_P:\mathbb{F}^n\to\mathbb{F}^n$  $T_P\mathbf{x}\defeq P\mathbf{x}$.\\
Set $\mathcal{U}\defeq\operatorname{range}(T_P)=\operatorname{col}(P)$. Then $T_P^2=T_P=T_P^*$ by $P^2=P=P^*$ and Theorem \ref{thm:matrix-of-adjoint-is-conjugate-transpose}.\\
Hence $T_P=P_{\mathcal{U}}$ by Theorem \ref{thm:orthogonal-projection-self-adjoint-idempotent}. Therefore $(T_P|_{\mathcal{U}},T_P|_{\mathcal{U}^{\perp}})=(I_{\mathcal{U}},0)$ by Proposition \ref{prop:orthogonal-projection-properties}(ii)--(iii).
\item \hlt{Adapted Basis}: Choose an orthonormal basis $(\mathbf{u}_1,\ldots,\mathbf{u}_r)$ of $\mathcal{U}$ by Corollary \ref{cor:finite-dimensional-inner-product-space-has-orthonormal-basis}.\\
Extend it to an orthonormal basis $(\mathbf{u}_1,\ldots,\mathbf{u}_n)$ of $\mathbb{F}^n$ by Corollary \ref{cor:orthonormal-list-extends-to-basis}.\\
Set $U\defeq[\mathbf{u}_1\ \cdots\ \mathbf{u}_n]$. Then $U^*U=I_n$ by Theorem \ref{thm:unitary-matrix-equivalent-conditions}(i)--(ii).
\item \hlt{Block Form}: If $j\in J_r$, then $U^*PU\mathbf{e}_j=U^*P\mathbf{u}_j=U^*\mathbf{u}_j=\mathbf{e}_j$. If $j\in J_n\setminus J_r$, then $U^*PU\mathbf{e}_j=\mathbf{0}$.\\
Thus $U^*PU=I_r\oplus0_{n-r}$ by equality on the standard basis.
\item \hlt{Similarity}: The identity $UU^*=I_n$ holds by Theorem \ref{thm:unitary-matrix-equivalent-conditions}(i),(iii).\\
Multiplication by $U$ on the left and $U^*$ on the right establishes $P=U(I_r\oplus0_{n-r})U^*$.
\end{enumerate}
\hlt{Structural Commentary}: The induced projection supplies an adapted orthonormal basis whose matrix action is block diagonal.\\
\hlt{Key Insight}: Projection is coordinate truncation in range--kernel coordinates.
\end{proof}
\end{theorem}

Figure \ref{fig:orthogonal-projection-unitary-block} depicts the range--kernel coordinates in Theorem \ref{thm:orthogonal-projection-matrix-unitary-coordinate-form}.
\begin{illustrationfigure}[H]
  \centering
  \begin{tikzpicture}[
  x=0.88cm,
  y=0.88cm,
  >=Stealth,
  every node/.style={font=\small},
  sub/.style={draw=blue!60!black, fill=blue!7, rounded corners=3pt, minimum width=2.07cm, minimum height=0.64cm, align=center},
  mat/.style={draw=gray!60, fill=gray!5, rounded corners=2pt, minimum width=1.90cm, minimum height=1.58cm},
  arrow/.style={->, line width=0.85pt, blue!70!black}
]
  \node[sub] (C) at (0,0.72) {$\operatorname{col}(P)$};
  \node[sub, fill=red!6, draw=red!55!black] (N) at (0,-0.72) {$\operatorname{col}(P)^\perp$};

  \node[mat] (B) at (4.1,0) {};
  \draw[gray!55] ($(B.center)+(-0.86,0)$) -- ($(B.center)+(0.86,0)$);
  \draw[gray!55] ($(B.center)+(0,-0.68)$) -- ($(B.center)+(0,0.68)$);
  \node[blue!70!black] at ($(B.center)+(-0.43,0.43)$) {$I_r$};
  \node[red!60!black] at ($(B.center)+(0.43,-0.43)$) {$0_{n-r}$};
  \node[gray!65] at ($(B.center)+(0.43,0.43)$) {$0$};
  \node[gray!65] at ($(B.center)+(-0.43,-0.43)$) {$0$};

  \draw[arrow] (C.east) -- ($(B.west)+(0,0.43)$);
  \draw[arrow, red!60!black] (N.east) -- ($(B.west)+(0,-0.43)$);
\end{tikzpicture}

  \caption{Unitary coordinates make an orthogonal projection the identity on $\operatorname{col}(P)$ and zero on $\operatorname{col}(P)^\perp$.}
  \label{fig:orthogonal-projection-unitary-block}
\end{illustrationfigure}

\begin{proposition} \label{prop:algebra-of-unitary-matrices} \hlt{Algebra of Unitary Matrices}\\
Let $m,n\geq1$, let $U,V\in\operatorname{Mat}(n,n,\mathbb{F})$, and let $W\in\operatorname{Mat}(m,m,\mathbb{F})$. Let $V\oplus W$ denote the block diagonal matrix with blocks $V$ and $W$.
\begin{enumerate}[label=(\roman*)]
\setlength{\itemsep}{0pt}
\item $[U^*U=I_n\land V^*V=I_n]\Longrightarrow(UV)^*(UV)=I_n$;
\item $[(V\oplus W)^*(V\oplus W)=I_{n+m}]\Longleftrightarrow[V^*V=I_n\land W^*W=I_m]$.
\end{enumerate}
\begin{proof}
\begin{enumerate}[label=(\roman*)]
\setlength{\itemsep}{0pt}
\item[$(i)$] The identity $(UV)^*(UV)=V^*U^*UV=V^*V=I_n$ holds by the hypotheses.
\item[$(ii)$] ($\Rightarrow$) If $(V\oplus W)^*(V\oplus W)=I_{n+m}$, then $V^*V\oplus W^*W=I_n\oplus I_m$.\\
Hence $V^*V=I_n$ and $W^*W=I_m$ by equality of the diagonal blocks.
\item[$(ii)$] ($\Leftarrow$) If $V^*V=I_n$ and $W^*W=I_m$, then $(V\oplus W)^*(V\oplus W)=I_n\oplus I_m=I_{n+m}$.
\end{enumerate}
\hlt{Structural Commentary}: Product reversal proves closure under multiplication; blockwise adjoints reduce direct sums componentwise.\\
\hlt{Key Insight}: Unitary identities multiply and split blockwise.
\end{proof}
\end{proposition}

\begin{theorem} \label{thm:extend-isometric-columns-to-unitary-matrix} \hlt{Extending Isometric Columns to a Unitary Matrix}\\
Let $m,n\geq1$ and let $X\in\operatorname{Mat}(m,n,\mathbb{F})$ satisfy $X^*X=I_n$. Then:
\begin{enumerate}[label=(\roman*)]
\setlength{\itemsep}{0pt}
\item $n\leq m$;
\item $(\exists\,X'\in\operatorname{Mat}(m,m-n,\mathbb{F}))\,[X\ X']^*[X\ X']=I_m$.
\end{enumerate}
\begin{proof}
Write $X=[\mathbf{x}_1\ \cdots\ \mathbf{x}_n]$. For $i,j\in J_n$, $(X^*X)_{ij}=\innerprod[\mathbb{F}^m]{\mathbf{x}_j,\mathbf{x}_i}=\delta_{ij}$ by hypothesis. Hence $(\mathbf{x}_1,\ldots,\mathbf{x}_n)$ is orthonormal in $\mathbb{F}^m$.
\begin{enumerate}[label=(\roman*)]
\setlength{\itemsep}{0pt}
\item The list is linearly independent by Corollary \ref{cor:orthonormal-list-linearly-independent}, so $n\leq m$ by Theorem \ref{thm:finitedimlenofvec}.
\item Extend the list to an orthonormal basis $(\mathbf{x}_1,\ldots,\mathbf{x}_m)$ of $\mathbb{F}^m$ by Corollary \ref{cor:orthonormal-list-extends-to-basis}.\\
Set $X'\defeq[\mathbf{x}_{n+1}\ \cdots\ \mathbf{x}_m]$. Then $[X\ X']^*[X\ X']=I_m$ by Theorem \ref{thm:unitary-matrix-equivalent-conditions}(i)--(ii).
\end{enumerate}
\hlt{Structural Commentary}: The Gram identity turns the columns into an orthonormal list, then basis extension completes it.\\
\hlt{Key Insight}: A partial orthonormal frame extends to a unitary frame.
\end{proof}
\end{theorem}
